\section{Ciclos medios por acceso a memoria}
\label{sec:cpi}

Sea $T$ el coste medio (en ciclos) de un acceso a memoria del MIPS. Distinguimos las
siguientes clases de eventos, con sus probabilidades $p_x$ (que suman 1) y costes
$C_x$:

\begin{table}[htbp]
\centering
\renewcommand{\arraystretch}{1.25}
\begin{tabular}{l l l}
\toprule
\textbf{Evento} & \textbf{Coste $C_x$} & \textbf{Notas} \\
\midrule
Hit lectura/escritura            & $1$ & resuelto en \estado{Inicio} en un ciclo \\
Miss limpio (lectura cacheable)  & $T_\textit{arb} + 1 + \mathrm{CrB(MD)}$
                                 & \estado{Arbitro}+\estado{ADDR}+\estado{Fetch} \\
Miss sucio (lectura cacheable)   & $T_\textit{arb} + 1 + \mathrm{CwB(MD)} + 1 + \mathrm{CrB(MD)}$
                                 & a\~nade copy-back y nueva fase ADDR \\
Write miss (write-around)        & $T_\textit{arb} + 1 + \mathrm{CwW(MD)}$
                                 & ni sube bloque ni incrementa \sig{cont\_w} \\
Acceso a MD Scratch (lw/sw)      & $T_\textit{arb} + 1 + \mathrm{CrW/CwW(MDscratch)}$ & non-cacheable \\
Acceso a IO\_REG                 & $\mathrm{CrW/CwW(IO\_REG)}$
                                 & no usa el bus \\
Lectura a \sig{Addr\_Error\_Reg} & $1$ & se sirve en \estado{Inicio} \\
Acceso err\'oneo                 & $T_\textit{arb} + 1 + 1$
                                 & \estado{Arbitro}+\estado{ADDR}+vuelta a \estado{Inicio} \\
\bottomrule
\end{tabular}
\caption{Coste por evento (\textit{$T_\textit{arb}=1.5$ medio}, $+1$ por la fase
\estado{ADDR}).}
\label{tab:costes}
\end{table}

\paragraph{F\'ormula del coste medio.} Sea $p_h$ la probabilidad de acierto,
$p_{m,c}$ la de miss limpio, $p_{m,d}$ la de miss sucio, $p_{wa}$ la de write
miss, $p_\textit{scr}$ la de acceso a scratch y $p_\textit{io}$ la de acceso
a IO. Entonces:
\begin{align}
T &= p_h\cdot 1
   + p_{m,c}\bigl(T_\textit{arb} + 1 + \mathrm{CrB(MD)}\bigr) \nonumber \\
  &\quad + p_{m,d}\bigl(T_\textit{arb} + 1 + \mathrm{CwB(MD)} + 1 + \mathrm{CrB(MD)}\bigr) \nonumber \\
  &\quad + p_{wa}\bigl(T_\textit{arb} + 1 + \mathrm{CwW(MD)}\bigr) \nonumber \\
  &\quad + p_\textit{scr}\bigl(T_\textit{arb} + 1 + C_\textit{scr}\bigr)
   + p_\textit{io}\cdot C_\textit{io}.
\label{eq:cpi}
\end{align}

\paragraph{Estimaci\'on para el test de integraci\'on.}
A partir de los contadores observados al final de \texttt{tb\_integracion.asm}
(\dato{cont\_m=13, cont\_w=4, cont\_r=20, cont\_cb=4} esperados; verificar tras
simular) y los accesos no contados (scratch e IO), el reparto de eventos es:
\begin{itemize}[leftmargin=1.2em]
  \item Aciertos cacheable = $\mathrm{cont\_r}+\mathrm{cont\_w} - \text{accesos sobre miss}$ \dato{TBD}.
  \item Misses limpios = $\mathrm{cont\_m}-\mathrm{cont\_cb}-N_\textit{wa}$ = \dato{TBD}.
  \item Misses sucios = $\mathrm{cont\_cb}=4$.
  \item Write-around = $N_\textit{wa}=1$ (la sw a \texttt{0x80} en F5).
\end{itemize}

Sustituyendo los valores de la tabla~\ref{tab:latencias} en la
ecuaci\'on~\eqref{eq:cpi} se obtiene $T = \dato{TBD}$ ciclos por acceso.
